\paragraph{Poisson--Cauchy 核在首个非消失矩阶上的显式常数链}
在 Poisson 粗粒化框架下，\(L^1\) 几何、相对熵主项与耗散率可统一归约为同一个“首个非消失中心矩阶”参数。以下给出该链条的闭式常数表达、分层刚性与缺陷测度口径下的可审计分类。

\begin{theorem}[标准 Poisson 核高阶导数的 \texorpdfstring{$L^1$}{L1} 范数与 Fisher 型积分闭式]\label{thm:xi-poisson-kernel-derivative-l1-fisher-closed}
\leanverified{paper\_xi\_poisson\_kernel\_derivative\_l1\_fisher\_closed}
令
\[
\mathsf P_t(x):=\frac1\pi\frac{t}{x^2+t^2},\qquad t>0,
\]
并记 \(\mathsf P_1^{(m)}:=\partial_x^m\mathsf P_1\)（\(m\in\mathbb N\)）。则对任意整数 \(m\ge 1\) 有
\begin{align}
\|\mathsf P_1^{(m)}\|_{L^1(\RR)}
&=\frac{2(m-1)!}{\pi}\sum_{k=1}^{m}\sin^{m+1}\!\Bigl(\frac{k\pi}{m+1}\Bigr),\label{eq:xi-poisson-kernel-derivative-l1-closed}\\
\int_{\RR}\frac{\bigl(\mathsf P_1^{(m)}(x)\bigr)^2}{\mathsf P_1(x)}\,\dd x
&=\frac{m(2m)!}{(2m-1)\,2^{2m}}.\label{eq:xi-poisson-kernel-derivative-fisher-closed}
\end{align}
\end{theorem}
\begin{proof}
设 \(f(x):=(1+x^2)^{-1}\)，则 \(\mathsf P_1(x)=\pi^{-1}f(x)\)。由分式分解
\[
f(x)=\frac{1}{2i}\Bigl((x-i)^{-1}-(x+i)^{-1}\Bigr),
\]
可得
\[
f^{(m)}(x)
=\frac{(-1)^m m!}{2i}\Bigl((x-i)^{-m-1}-(x+i)^{-m-1}\Bigr)
=(-1)^{m+1}m!\,\Im\!\bigl((x+i)^{-m-1}\bigr).
\]
令 \(x=\cot\theta\)（\(\theta\in(0,\pi)\)），则
\[
x+i=\frac{e^{i\theta}}{\sin\theta},
\]
从而
\begin{equation}\label{eq:xi-poisson-kernel-derivative-trig-form}
\mathsf P_1^{(m)}(\cot\theta)
=\frac{(-1)^m m!}{\pi}\sin^{m+1}\!\theta\,\sin\!\bigl((m+1)\theta\bigr).
\end{equation}

先证 \eqref{eq:xi-poisson-kernel-derivative-fisher-closed}。由
\[
\dd x=\csc^2\theta\,\dd\theta,\qquad
\mathsf P_1(\cot\theta)=\frac1\pi\sin^2\theta,
\]
代入 \eqref{eq:xi-poisson-kernel-derivative-trig-form} 得
\[
\int_{\RR}\frac{(\mathsf P_1^{(m)})^2}{\mathsf P_1}\,\dd x
=\frac{m!^2}{\pi}
\int_0^\pi
\sin^{2m-2}\!\theta\,
\sin^2\!\bigl((m+1)\theta\bigr)\,\dd\theta.
\]
写
\(
\sin^2((m+1)\theta)=\frac12\bigl(1-\cos(2(m+1)\theta)\bigr)
\)。
由于 \(\sin^{2m-2}\theta\) 的 Fourier 频率仅到 \(2(m-1)\)，与 \(\cos(2(m+1)\theta)\) 在 \([0,\pi]\) 上正交，故
\[
\int_0^\pi\sin^{2m-2}\!\theta\,\sin^2((m+1)\theta)\,\dd\theta
=\frac12\int_0^\pi\sin^{2m-2}\!\theta\,\dd\theta.
\]
再用经典积分
\[
\int_0^\pi\sin^{2(m-1)}\!\theta\,\dd\theta
=\pi\,\frac{(2m-2)!}{2^{2m-2}((m-1)!)^2},
\]
得到
\[
\int_{\RR}\frac{(\mathsf P_1^{(m)})^2}{\mathsf P_1}\,\dd x
=\frac{m!^2(2m-2)!}{2^{2m-1}((m-1)!)^2}
=\frac{m(2m)!}{(2m-1)\,2^{2m}}.
\]

再证 \eqref{eq:xi-poisson-kernel-derivative-l1-closed}。由
\eqref{eq:xi-poisson-kernel-derivative-trig-form}，
\(\mathsf P_1^{(m)}\) 的零点对应
\(
\sin((m+1)\theta)=0
\)，即
\[
\theta_k=\frac{k\pi}{m+1}\ (k=1,\dots,m),\qquad
x_k=\cot\theta_k.
\]
这些零点均为单零点，且
\(
\lim_{x\to\pm\infty}\mathsf P_1^{(m-1)}(x)=0
\)。
故总变差恒等式给出
\begin{equation}\label{eq:xi-poisson-kernel-derivative-tv}
\|\mathsf P_1^{(m)}\|_{L^1(\RR)}
=\mathrm{TV}\!\bigl(\mathsf P_1^{(m-1)}\bigr)
=2\sum_{k=1}^{m}\bigl|\mathsf P_1^{(m-1)}(x_k)\bigr|.
\end{equation}
将 \eqref{eq:xi-poisson-kernel-derivative-trig-form} 中 \(m\) 换为 \(m-1\)，并取 \(\theta=\theta_k\)，有
\[
\sin(m\theta_k)=\sin(k\pi-\theta_k)=(-1)^{k+1}\sin\theta_k,
\]
从而
\[
\bigl|\mathsf P_1^{(m-1)}(x_k)\bigr|
=\frac{(m-1)!}{\pi}\sin^{m+1}\!\theta_k.
\]
代回 \eqref{eq:xi-poisson-kernel-derivative-tv} 即得 \eqref{eq:xi-poisson-kernel-derivative-l1-closed}。证毕。
\end{proof}

\begin{theorem}[Poisson--Cauchy 粗粒化在首个非消失矩阶上的 \texorpdfstring{$L^1$}{L1}/KL 锐渐近]\label{thm:xi-poisson-first-nonvanishing-moment-tv-kl-sharp}
设 \(\nu\) 为 \(\RR\) 上概率测度，\(c\in\RR\)，并定义
\[
g_t(x):=\mathsf P_t(x-c),\qquad
h_t(x):=\int_{\RR}\mathsf P_t(x-\gamma)\,\dd\nu(\gamma).
\]
对 \(m\ge 1\)，定义中心矩
\[
\kappa_k:=\int_{\RR}(\gamma-c)^k\,\dd\nu(\gamma)\qquad(k\in\mathbb N).
\]
若满足
\[
\int_{\RR}|\gamma-c|^{m+1}\,\dd\nu(\gamma)<\infty,\qquad
\kappa_1=\cdots=\kappa_{m-1}=0,\qquad
\kappa_m\neq 0,
\]
则当 \(t\to\infty\) 有
\begin{align}
t^m\|h_t-g_t\|_{L^1(\RR)}
&\longrightarrow
\frac{|\kappa_m|}{m!}\,\|\mathsf P_1^{(m)}\|_{L^1(\RR)}\notag\\
&=
\frac{2|\kappa_m|}{m\pi}
\sum_{k=1}^{m}\sin^{m+1}\!\Bigl(\frac{k\pi}{m+1}\Bigr),\label{eq:xi-poisson-first-nonvanishing-moment-tv-sharp}\\
t^{2m}D_{\mathrm{KL}}(h_t\mid g_t)
&\longrightarrow
\frac{\kappa_m^2}{2(m!)^2}
\int_{\RR}\frac{\bigl(\mathsf P_1^{(m)}(x)\bigr)^2}{\mathsf P_1(x)}\,\dd x\notag\\
&=
\frac{\kappa_m^2}{(m!)^2}\cdot
\frac{m(2m)!}{(2m-1)\,2^{2m+1}}.\label{eq:xi-poisson-first-nonvanishing-moment-kl-sharp}
\end{align}
\end{theorem}
\begin{proof}
记 \(\Phi_t(x,u):=\mathsf P_t(x-u)\)。对固定 \(x\)，\(u\mapsto \Phi_t(x,u)\) 为 \(C^\infty\)，且
\[
\partial_u^k\Phi_t(x,u)=(-1)^k\partial_x^k\mathsf P_t(x-u).
\]
在 \(u=c\) 处对 \(u=\gamma\) 作 \(m\) 阶 Taylor 展开（积分余项形式）并对 \(\nu\) 积分，得
\[
h_t(x)-g_t(x)=\frac{\kappa_m}{m!}\,\partial_x^m g_t(x)+\varepsilon_t(x),
\]
其中
\[
\varepsilon_t(x)
=\int_{\RR}R_{m+1}(x,\gamma)\,\dd\nu(\gamma),
\]
且
\[
\|\varepsilon_t\|_{L^1}
\le
\frac{\int |\gamma-c|^{m+1}\dd\nu(\gamma)}{(m+1)!}
\cdot
\|\partial_x^{m+1}\mathsf P_t\|_{L^1}
=O(t^{-m-1}).
\]
这里使用了尺度关系
\[
\partial_x^r\mathsf P_t(x)=t^{-r-1}\mathsf P_1^{(r)}(x/t),
\qquad
\|\partial_x^r\mathsf P_t\|_{L^1}=t^{-r}\|\mathsf P_1^{(r)}\|_{L^1}.
\]
故
\[
\|h_t-g_t\|_{L^1}
=\frac{|\kappa_m|}{m!}\,t^{-m}\|\mathsf P_1^{(m)}\|_{L^1}+o(t^{-m}),
\]
得到 \eqref{eq:xi-poisson-first-nonvanishing-moment-tv-sharp}。

再看 KL 项。令
\(
\delta_t:=h_t/g_t-1
\)。
由上式可写
\[
\delta_t(x)
=\frac{\kappa_m}{m!}\frac{\partial_x^m g_t(x)}{g_t(x)}+\frac{\varepsilon_t(x)}{g_t(x)}.
\]
写 \(y=(x-c)/t\)，则
\[
\frac{\partial_x^m g_t(x)}{g_t(x)}
=t^{-m}\frac{\mathsf P_1^{(m)}(y)}{\mathsf P_1(y)}.
\]
并且 \(\mathsf P_1^{(m)}/\mathsf P_1\) 为有界有理函数。结合余项估计，可得
\[
\delta_t
=\frac{\kappa_m}{m!}\,t^{-m}\frac{\mathsf P_1^{(m)}(y)}{\mathsf P_1(y)}
+o_{L^2(g_t)}(t^{-m}),
\]
且 \(\|\delta_t\|_\infty\to 0\)。
由展开
\(
(1+u)\log(1+u)=u+\frac12u^2+o(u^2)
\)
与
\(
\int g_t\delta_t\,\dd x=\int(h_t-g_t)\dd x=0
\)，得
\[
D_{\mathrm{KL}}(h_t\mid g_t)
=\frac12\int_{\RR}g_t\,\delta_t^2\,\dd x+o(t^{-2m}).
\]
乘以 \(t^{2m}\)，再作 \(x=c+ty\) 代换，即
\[
t^{2m}D_{\mathrm{KL}}(h_t\mid g_t)
\to
\frac{\kappa_m^2}{2(m!)^2}
\int_{\RR}\frac{\bigl(\mathsf P_1^{(m)}(y)\bigr)^2}{\mathsf P_1(y)}\,\dd y.
\]
最后代入定理
\ref{thm:xi-poisson-kernel-derivative-l1-fisher-closed}
的 \eqref{eq:xi-poisson-kernel-derivative-fisher-closed} 得
\eqref{eq:xi-poisson-first-nonvanishing-moment-kl-sharp}。证毕。
\end{proof}

\leanverified{paper\_xi\_poisson\_golden\_chain\_entropy\_slope\_m}
\begin{corollary}[KL 衰减阶在黄金尺度链上的可审计反演]\label{cor:xi-poisson-golden-chain-entropy-slope-m}
在定理 \ref{thm:xi-poisson-first-nonvanishing-moment-tv-kl-sharp} 的假设下，令 \(t_n:=\varphi^n\) 并定义熵斜率
\[
S_n:=-\frac1n\log D_{\mathrm{KL}}(h_{t_n}\mid g_{t_n}).
\]
则极限存在并满足
\[
\lim_{n\to\infty}S_n=2m\log\varphi,
\qquad
m=\lim_{n\to\infty}\frac{S_n}{2\log\varphi}.
\]
\end{corollary}
\begin{proof}
由 \eqref{eq:xi-poisson-first-nonvanishing-moment-kl-sharp}，存在常数 \(C>0\) 使
\(
D_{\mathrm{KL}}(h_t\mid g_t)=C\,t^{-2m}(1+o(1))
\)（\(t\to\infty\)）。
代入 \(t=t_n=\varphi^n\) 得
\[
-\log D_{\mathrm{KL}}(h_{t_n}\mid g_{t_n})
=2m\,n\log\varphi-\log C+o(1),
\]
两边除以 \(n\) 并令 \(n\to\infty\) 即得所示极限与反演式。证毕。
\end{proof}

\begin{corollary}[KL 与 \texorpdfstring{$L^1$}{L1} 平方等价的渐近常数]\label{cor:xi-poisson-kl-l1-square-equivalence-first-moment}
在定理 \ref{thm:xi-poisson-first-nonvanishing-moment-tv-kl-sharp} 的假设下，存在显式常数
\[
\mathfrak C_m
:=
\frac{\displaystyle\frac{m(2m)!}{(2m-1)\,2^{2m+1}}}
{\bigl\|\mathsf P_1^{(m)}\bigr\|_{L^1(\RR)}^2}
\]
使得当 \(t\to\infty\) 时成立二次等价
\[
D_{\mathrm{KL}}(h_t\mid g_t)
=\mathfrak C_m\,\|h_t-g_t\|_{L^1(\RR)}^{\,2}\,\bigl(1+o(1)\bigr).
\]
\end{corollary}
\begin{proof}
将 \eqref{eq:xi-poisson-first-nonvanishing-moment-tv-sharp} 与
\eqref{eq:xi-poisson-first-nonvanishing-moment-kl-sharp} 的极限常数相除，\(\kappa_m^2/(m!)^2\) 抵消，得到
\[
\frac{D_{\mathrm{KL}}(h_t\mid g_t)}{\|h_t-g_t\|_{L^1(\RR)}^2}\to \mathfrak C_m.
\]
即为所述二次等价。证毕。
\end{proof}

\begin{proposition}[二阶可微 \texorpdfstring{$f$}{f}-散度首项由 \texorpdfstring{$f''(1)$}{f''(1)} 与首个非消失矩阶完全决定]\label{prop:xi-poisson-first-nonvanishing-moment-fdiv-sharp}
\leanverified{paper\_xi\_poisson\_first\_nonvanishing\_moment\_fdiv\_sharp}
在定理 \ref{thm:xi-poisson-first-nonvanishing-moment-tv-kl-sharp} 的假设下，
令 \(f\) 为凸函数，且在 \(1\) 邻域二阶可微，满足
\(
f(1)=f'(1)=0
\)。
定义
\[
D_f(h_t\mid g_t):=\int_{\RR}g_t(x)\,
f\!\left(\frac{h_t(x)}{g_t(x)}\right)\dd x.
\]
则
\begin{equation}\label{eq:xi-poisson-first-nonvanishing-moment-fdiv-sharp}
t^{2m}D_f(h_t\mid g_t)
\longrightarrow
\frac{f''(1)}{2}\cdot\frac{\kappa_m^2}{(m!)^2}\cdot
\frac{m(2m)!}{(2m-1)\,2^{2m}}.
\end{equation}
\end{proposition}
\begin{proof}
由定理 \ref{thm:xi-poisson-first-nonvanishing-moment-tv-kl-sharp} 证明中的记号，
\(\delta_t:=h_t/g_t-1\to 0\)，且
\[
\delta_t
=\frac{\kappa_m}{m!}\,t^{-m}\frac{\mathsf P_1^{(m)}(y)}{\mathsf P_1(y)}
+o_{L^2(g_t)}(t^{-m}).
\]
由
\(
f(1+u)=\frac{f''(1)}2u^2+o(u^2)
\)
与 \(\int g_t\delta_t\,\dd x=0\)，可得
\[
D_f(h_t\mid g_t)
=\frac{f''(1)}2\int g_t\delta_t^2\,\dd x+o(t^{-2m}).
\]
再用
\[
\int g_t\delta_t^2\,\dd x
\sim
\frac{\kappa_m^2}{(m!)^2}\,t^{-2m}
\int_{\RR}\frac{(\mathsf P_1^{(m)}(y))^2}{\mathsf P_1(y)}\,\dd y
\]
并代入 \eqref{eq:xi-poisson-kernel-derivative-fisher-closed} 即得
\eqref{eq:xi-poisson-first-nonvanishing-moment-fdiv-sharp}。证毕。
\end{proof}

\begin{corollary}[KL 耗散率 \texorpdfstring{$I_1$}{I1} 的阶数与显式常数]\label{cor:xi-poisson-first-nonvanishing-moment-i1-sharp}
在定理 \ref{thm:xi-poisson-first-nonvanishing-moment-tv-kl-sharp} 的设定下，令
\[
I_1(h_t\mid g_t):=-\frac{\dd}{\dd t}D_{\mathrm{KL}}(h_t\mid g_t).
\]
则
\begin{equation}\label{eq:xi-poisson-first-nonvanishing-moment-i1-sharp}
t^{2m+1}I_1(h_t\mid g_t)
\longrightarrow
\frac{\kappa_m^2}{(m!)^2}\cdot
\frac{m^2(2m)!}{(2m-1)\,2^{2m}}.
\end{equation}
\end{corollary}
\begin{proof}
由 \eqref{eq:xi-poisson-first-nonvanishing-moment-kl-sharp}，
\[
D_{\mathrm{KL}}(h_t\mid g_t)=C\,t^{-2m}+o(t^{-2m}),
\quad
C:=\frac{\kappa_m^2}{(m!)^2}\cdot\frac{m(2m)!}{(2m-1)\,2^{2m+1}}.
\]
对 \(t\) 求导并用定义
\(
I_1=-\frac{\dd}{\dd t}D_{\mathrm{KL}}
\)
即得
\[
I_1(h_t\mid g_t)=2mC\,t^{-2m-1}+o(t^{-2m-1}),
\]
从而得到 \eqref{eq:xi-poisson-first-nonvanishing-moment-i1-sharp}。证毕。
\end{proof}

\begin{proposition}[KL 衰减阶对矩消失的必要条件]\label{prop:xi-poisson-kl-rate-implies-moment-cancellation}
\leanverified{paper\_xi\_poisson\_kl\_rate\_implies\_moment\_cancellation}
设 \(\nu\) 具有所有阶有限矩，固定 \(c\in\RR\)，并记
\[
h_t=\mathsf P_t*\nu,\qquad
g_t=\mathsf P_t(\cdot-c).
\]
对任意 \(m\ge 1\)，若
\[
D_{\mathrm{KL}}(h_t\mid g_t)=o(t^{-2m})\qquad(t\to\infty),
\]
则必有
\[
\int_{\RR}(\gamma-c)^k\,\dd\nu(\gamma)=0,\qquad k=1,\dots,m.
\]
特别地，若对所有 \(m\) 都有
\(
D_{\mathrm{KL}}(h_t\mid g_t)=o(t^{-2m})
\)，则 \(\nu=\delta_c\)。
\end{proposition}
\begin{proof}
设
\[
m_\ast:=\min\{k\ge 1:\ \kappa_k\neq 0\},
\qquad
\kappa_k:=\int_{\RR}(\gamma-c)^k\,\dd\nu(\gamma).
\]
若 \(m_\ast\) 存在，则由定理
\ref{thm:xi-poisson-first-nonvanishing-moment-tv-kl-sharp}
（取 \(m=m_\ast\)）有
\[
D_{\mathrm{KL}}(h_t\mid g_t)\sim
\frac{\kappa_{m_\ast}^2}{(m_\ast!)^2}
\frac{m_\ast(2m_\ast)!}{(2m_\ast-1)\,2^{2m_\ast+1}}
t^{-2m_\ast},
\]
与 \(o(t^{-2m_\ast})\) 矛盾，故 \(\kappa_1=\cdots=\kappa_m=0\)。
若 \(m_\ast\) 不存在，则全部中心矩为零，尤其
\[
\int_{\RR}|\gamma-c|^{2r}\,\dd\nu(\gamma)=0\qquad(\forall r\ge 1),
\]
从而 \(\gamma=c\) \(\nu\)-a.s.，即 \(\nu=\delta_c\)。证毕。
\end{proof}

\begin{theorem}[缺陷 Poisson 粗粒化的多极指纹分类]\label{thm:xi-poisson-defect-measure-multipole-fingerprint-classification}
在缺陷测度构造
\eqref{eq:xi-offcritical-defect-measure-nu}--\eqref{eq:xi-offcritical-comoving-poisson-field}
下，定义归一化高度测度
\[
\tilde\nu:=\frac1{M_0}\sum_j m_j\delta_{\gamma_j},
\]
以及
\[
\tilde h_t(x):=\int_{\RR}\mathsf P_t(x-\gamma)\,\dd\tilde\nu(\gamma),\qquad
\tilde g_t(x):=\mathsf P_t(x-\bar\gamma),
\]
其中 \(\bar\gamma:=\int \gamma\,\dd\tilde\nu(\gamma)\)。
再定义
\[
\tilde\kappa_k:=\int_{\RR}(\gamma-\bar\gamma)^k\,\dd\tilde\nu(\gamma).
\]
若存在最小整数 \(m_\ast\ge 1\) 使
\[
\tilde\kappa_{m_\ast}\neq 0,\qquad
\int_{\RR}|\gamma-\bar\gamma|^{m_\ast+1}\,\dd\tilde\nu(\gamma)<\infty,
\]
则当 \(t\to\infty\)：
\begin{enumerate}
\item \(\|\tilde h_t-\tilde g_t\|_{L^1}\) 的主导衰减阶为 \(t^{-m_\ast}\)，且主项常数由
\eqref{eq:xi-poisson-first-nonvanishing-moment-tv-sharp}
（以 \(\kappa_{m_\ast}\) 替换为 \(\tilde\kappa_{m_\ast}\)）给出；
\item \(D_{\mathrm{KL}}(\tilde h_t\mid \tilde g_t)\) 的主导衰减阶为 \(t^{-2m_\ast}\)，且主项常数由
\eqref{eq:xi-poisson-first-nonvanishing-moment-kl-sharp}
（同样替换 \(\kappa_{m_\ast}\)）给出；
\item 若观测到
\(
D_{\mathrm{KL}}(\tilde h_t\mid \tilde g_t)=o(t^{-2m})
\)，则必有
\(
\tilde\kappa_1=\cdots=\tilde\kappa_m=0
\)。
\end{enumerate}
\end{theorem}
\begin{proof}
取定理 \ref{thm:xi-poisson-first-nonvanishing-moment-tv-kl-sharp} 与命题
\ref{prop:xi-poisson-kl-rate-implies-moment-cancellation}
中的 \(\nu=\tilde\nu\)、\(c=\bar\gamma\)，即可逐条得到结论。证毕。
\end{proof}

\endinput
